% ============================================================
% Chapter 6 — Firmware Architecture & Software Design
% Source: PSYCON Engineering Design Document, Vol. 1, Ch. 6
% ============================================================
\section{Firmware Architecture \& Software Design}
\label{sec:firmware-architecture}

\subsection{Firmware Overview}
\label{sec:firmware-overview}

PSYCON uses two independent ESP32 firmware projects, one per module.

\textbf{Wrist Module firmware} is responsible for: MAX30102 acquisition; MPU6050 acquisition; MCP9808 acquisition; GSR acquisition (via the ADS1115); timestamping; local preprocessing; BLE/Wi-Fi transmission; battery monitoring; and error logging.

\textbf{Audio Module firmware} is responsible for: INMP441 audio capture; TEMT6000 acquisition; audio buffering; voice preprocessing; BLE/Wi-Fi transmission; battery monitoring; and error logging.

\subsection{Software Layers}
\label{sec:software-layers}

\Cref{fig:software-layers} shows the layered firmware architecture shared by both modules, from the application layer down to the physical ESP32 hardware.

\begin{figure}[H]
\centering
\begin{tikzpicture}[
    font=\small,
    node distance=0mm,
    layer/.style={draw, rounded corners=3pt, minimum width=6.4cm, align=center, line width=0.6pt},
    apps/.style={layer, fill=blue!8, minimum height=1.9cm},
    hal/.style={layer, fill=orange!10, minimum height=0.8cm},
    core/.style={layer, fill=gray!10, minimum height=0.8cm},
    hw/.style={layer, fill=green!8, minimum height=0.8cm}
]
\node[apps] (app) {
  \textbf{Application Layer}\\[2pt]
  \begin{tabular}{@{}l@{}}
  \textbullet\ Sensor Manager \\
  \textbullet\ Data Fusion \\
  \textbullet\ Communications \\
  \textbullet\ Battery Manager \\
  \textbullet\ Diagnostics
  \end{tabular}
};
\node[hal, below=5mm of app] (hal) {Hardware Abstraction Layer (HAL)};
\node[core, below=5mm of hal] (core) {ESP-IDF / Arduino Core};
\node[hw, below=5mm of core] (hw) {ESP32 Hardware};

\draw[-{Stealth[length=2mm]}, line width=0.6pt] (app.south) -- (hal.north);
\draw[-{Stealth[length=2mm]}, line width=0.6pt] (hal.south) -- (core.north);
\draw[-{Stealth[length=2mm]}, line width=0.6pt] (core.south) -- (hw.north);
\end{tikzpicture}
\caption{Layered firmware architecture: an application layer (sensor management, data fusion, communications, battery management, diagnostics) sits above a hardware abstraction layer, the ESP-IDF/Arduino core, and the ESP32 hardware itself.}
\label{fig:software-layers}
\end{figure}

\subsection{Wrist Module Tasks}
\label{sec:wrist-module-tasks}

The Wrist Module firmware is organized into eight cooperating tasks, summarized in \Cref{tab:wrist-tasks-summary} and illustrated in \Cref{fig:wrist-task-pipeline}.

\begin{table*}[t]
\centering
\caption{Wrist Module firmware task summary.}
\label{tab:wrist-tasks-summary}
\small
\begin{tabularx}{\textwidth}{@{}l l Y@{}}
\toprule
\textbf{Task} & \textbf{Trigger} & \textbf{Responsibilities} \\
\midrule
1. Sensor Initialization & Runs once & Initializes MAX30102, ADS1115, MPU6050, MCP9808 \\
2. Heart Rate Task & Periodic & Reads RED LED; reads IR LED; filters signal; stores buffer \\
3. Motion Task & Periodic & Reads accelerometer and gyroscope; computes movement and orientation \\
4. Temperature Task & Periodic & Reads MCP9808; stores temperature \\
5. GSR Task & Periodic & Reads ADS1115; converts ADC value to conductance; stores value \\
6. Packet Builder & Periodic & Assembles one packet from timestamp, heart, motion, temperature, GSR, and checksum \\
7. BLE/Wi-Fi Task & Periodic & Uploads packets; retries on failure \\
8. Battery Monitor & Periodic & Checks battery voltage; issues low-battery / shutdown warnings \\
\bottomrule
\end{tabularx}
\end{table*}

\begin{figure*}[t]
\centering
\begin{tikzpicture}[
    font=\scriptsize,
    stage/.style={draw, rounded corners=3pt, text width=1.6cm, minimum height=1.15cm, align=center, line width=0.6pt, fill=blue!6, inner sep=2pt},
    final/.style={stage, fill=green!8},
    arrow/.style={-{Stealth[length=2mm]}, line width=0.6pt},
    node distance=2mm
]
\node[stage] (t1) {Task 1\\Sensor Init};
\node[stage, right=of t1] (t2) {Task 2\\Heart Rate};
\node[stage, right=of t2] (t3) {Task 3\\Motion};
\node[stage, right=of t3] (t4) {Task 4\\Temperature};
\node[stage, right=of t4] (t5) {Task 5\\GSR};
\node[stage, right=of t5] (t6) {Task 6\\Packet Builder};
\node[stage, right=of t6] (t7) {Task 7\\BLE/Wi-Fi};
\node[final, right=of t7] (t8) {Task 8\\Battery Monitor};

\draw[arrow] (t1.east) -- (t2.west);
\draw[arrow] (t2.east) -- (t3.west);
\draw[arrow] (t3.east) -- (t4.west);
\draw[arrow] (t4.east) -- (t5.west);
\draw[arrow] (t5.east) -- (t6.west);
\draw[arrow] (t6.east) -- (t7.west);
\draw[arrow] (t7.east) -- (t8.west);
\end{tikzpicture}
\caption{Wrist Module task pipeline: one-time sensor initialization followed by periodic acquisition tasks (heart rate, motion, temperature, GSR), packet assembly, wireless upload, and battery monitoring.}
\label{fig:wrist-task-pipeline}
\end{figure*}

\subsection{Audio Tasks}
\label{sec:audio-tasks-ch6}

The Audio Module firmware runs four cooperating tasks: the Microphone Task, the Light Sensor Task, the Packet Builder, and the Battery Monitor / Communication task. \Cref{fig:audio-task-pipeline} shows the audio acquisition pipeline from initialization through transmission.

\begin{figure}[H]
\centering
\begin{tikzpicture}[
    font=\small,
    node distance=0mm,
    field/.style={draw, rounded corners=2pt, fill=violet!6, minimum width=3.6cm, minimum height=0.65cm, align=center, line width=0.5pt}
]
\node[field] (init) {Initialization};
\node[field, below=5mm of init] (i2s) {I\textsuperscript{2}S Configure};
\node[field, below=5mm of i2s] (dma) {DMA Buffers};
\node[field, below=5mm of dma] (rec) {Continuous Recording};
\node[field, below=5mm of rec] (frame) {Frame Buffer};
\node[field, below=5mm of frame] (tx) {Transmission};

\foreach \a/\b in {init/i2s, i2s/dma, dma/rec, rec/frame, frame/tx}
  \draw[-{Stealth[length=2mm]}, line width=0.6pt] (\a.south) -- (\b.north);
\end{tikzpicture}
\caption{Audio Module acquisition pipeline: initialization, I\textsuperscript{2}S configuration, DMA buffer setup, continuous recording, frame buffering, and transmission.}
\label{fig:audio-task-pipeline}
\end{figure}

\subsection{Data Flow}
\label{sec:firmware-data-flow}

\Cref{fig:firmware-data-flow} shows the end-to-end firmware data flow, from raw sensor reads through to server-side AI processing.

\begin{figure*}[t]
\centering
\begin{tikzpicture}[
    font=\scriptsize,
    stage/.style={draw, rounded corners=3pt, text width=1.6cm, minimum height=1.1cm, align=center, line width=0.6pt, fill=orange!8, inner sep=2pt},
    arrow/.style={-{Stealth[length=2mm]}, line width=0.6pt},
    node distance=2.5mm
]
\node[stage] (s1) {Sensors};
\node[stage, right=of s1] (s2) {Driver};
\node[stage, right=of s2] (s3) {Filter};
\node[stage, right=of s3] (s4) {Timestamp};
\node[stage, right=of s4] (s5) {Packet};
\node[stage, right=of s5] (s6) {Transmission};
\node[stage, right=of s6] (s7) {Server};
\node[stage, right=of s7] (s8) {AI};

\foreach \i/\j in {s1/s2, s2/s3, s3/s4, s4/s5, s5/s6, s6/s7, s7/s8}
  \draw[arrow] (\i.east) -- (\j.west);
\end{tikzpicture}
\caption{End-to-end firmware data flow, from sensor acquisition through driver-level reads, filtering, timestamping, packet assembly, wireless transmission, server ingestion, and AI processing.}
\label{fig:firmware-data-flow}
\end{figure*}

\subsection{Packet Structure}
\label{sec:packet-structure}

\Cref{tab:packet-physio} lists an example physiological packet structure (Wrist Module), and \Cref{tab:packet-audio} lists the audio packet structure (Audio Module).

\begin{table}[H]
\centering
\caption{Example physiological (Wrist Module) packet fields.}
\label{tab:packet-physio}
\small
\begin{tabularx}{\columnwidth}{@{}l Y@{}}
\toprule
\textbf{\#} & \textbf{Field} \\
\midrule
1 & Timestamp \\
2 & Heart rate \\
3 & SpO\textsubscript{2} \\
4 & Accelerometer X \\
5 & Accelerometer Y \\
6 & Accelerometer Z \\
7 & Gyroscope X \\
8 & Gyroscope Y \\
9 & Gyroscope Z \\
10 & Temperature \\
11 & GSR \\
12 & Battery \\
13 & Checksum \\
\bottomrule
\end{tabularx}
\end{table}

\begin{table}[H]
\centering
\caption{Audio packet fields.}
\label{tab:packet-audio}
\small
\begin{tabularx}{\columnwidth}{@{}l Y@{}}
\toprule
\textbf{\#} & \textbf{Field} \\
\midrule
1 & Timestamp \\
2 & PCM samples \\
3 & Battery \\
4 & Checksum \\
\bottomrule
\end{tabularx}
\end{table}

\subsection{Sampling Strategy}
\label{sec:sampling-strategy}

\Cref{tab:sampling-strategy} summarizes the acquisition regime for each signal. Actual sampling frequencies are to be finalized based on testing and downstream processing requirements.

\begin{table}[H]
\centering
\caption{Sampling strategy by signal.}
\label{tab:sampling-strategy}
\small
\begin{tabularx}{\columnwidth}{@{}Y l@{}}
\toprule
\textbf{Signal} & \textbf{Regime} \\
\midrule
Heart & Continuous \\
Motion & Continuous \\
Temperature & Periodic \\
Light & Periodic \\
Audio & Continuous \\
GSR & Continuous \\
\bottomrule
\end{tabularx}
\end{table}

\begin{requirementbox}
Actual sampling frequencies for each signal remain to be finalized based on empirical testing and the requirements of downstream processing.
\end{requirementbox}

\subsection{Error Handling}
\label{sec:error-handling-ch6}

\Cref{alg:error-handling} describes the generic sensor error-handling procedure used across both firmware projects. The system is designed so that a single failed sensor never causes a complete crash.

\begin{algorithm}[t]
\small
\caption{Generic Sensor Error Handling}
\label{alg:error-handling}
\KwData{Sensor read request}
\KwResult{Sensor value or logged failure, without halting the system}
Attempt sensor read\;
\If{timeout or communication error detected}{
  Retry read\;
  \If{retry fails}{
    Log error (sensor ID, error type, timestamp)\;
  }
}
Continue acquisition for all remaining sensors\;
\end{algorithm}

\subsection{Communication}
\label{sec:communication-ch6}

Every transmitted packet contains: a timestamp; a sensor ID; the payload; a CRC/checksum; and a sequence number. Because each packet carries a sequence number, missing packets can be detected at the receiving end.

\subsection{Logging}
\label{sec:logging-ch6}

The firmware maintains logs for: sensor initialization failures; I\textsuperscript{2}C communication errors; buffer overruns; Wi-Fi/BLE reconnects; battery warnings; brownout events; and memory allocation failures.

\subsection{Battery Management}
\label{sec:battery-management-ch6}

\Cref{fig:battery-state} shows the firmware battery-management state progression. Sudden power loss is avoided where possible by saving pending data before shutdown.

\begin{figure}[H]
\centering
\begin{tikzpicture}[
    font=\small,
    node distance=0mm,
    field/.style={draw, rounded corners=2pt, fill=red!6, minimum width=4.0cm, minimum height=0.65cm, align=center, line width=0.5pt}
]
\node[field, fill=green!8] (normal) {Normal};
\node[field, below=5mm of normal] (warn) {Battery Warning};
\node[field, below=5mm of warn] (save) {Save Pending Data};
\node[field, below=5mm of save] (shut) {Graceful Shutdown};

\foreach \a/\b in {normal/warn, warn/save, save/shut}
  \draw[-{Stealth[length=2mm]}, line width=0.6pt] (\a.south) -- (\b.north);
\end{tikzpicture}
\caption{Firmware battery-management state progression, from normal operation through a battery-warning state, saving of pending data, and graceful shutdown.}
\label{fig:battery-state}
\end{figure}

\subsection{Boot Sequence}
\label{sec:boot-sequence}

\Cref{fig:boot-sequence} shows the firmware boot sequence, from power-on to continuous acquisition.

\begin{figure}[H]
\centering
\begin{tikzpicture}[
    font=\small,
    node distance=0mm,
    field/.style={draw, rounded corners=2pt, fill=blue!6, minimum width=4.4cm, minimum height=0.6cm, align=center, line width=0.5pt}
]
\node[field] (pwr) {Power On};
\node[field, below=4mm of pwr] (boot) {ESP32 Boot};
\node[field, below=4mm of boot] (hwinit) {Initialize Hardware};
\node[field, below=4mm of hwinit] (sensinit) {Initialize Sensors};
\node[field, below=4mm of sensinit] (scan) {I\textsuperscript{2}C Scan};
\node[field, below=4mm of scan] (self) {Self-Test};
\node[field, below=4mm of self] (comm) {Communication Ready};
\node[field, below=4mm of comm, fill=green!8] (acq) {Continuous Acquisition};

\foreach \a/\b in {pwr/boot, boot/hwinit, hwinit/sensinit, sensinit/scan, scan/self, self/comm, comm/acq}
  \draw[-{Stealth[length=2mm]}, line width=0.6pt] (\a.south) -- (\b.north);
\end{tikzpicture}
\caption{Firmware boot sequence, from power-on through hardware and sensor initialization, an I\textsuperscript{2}C bus scan, self-test, and communication readiness, to continuous acquisition.}
\label{fig:boot-sequence}
\end{figure}

\subsection{Watchdog Recovery}
\label{sec:watchdog-recovery}

\begin{engnote}
If the firmware hangs, the watchdog resets the ESP32. The boot reason is logged, and data acquisition resumes automatically after restart.
\end{engnote}

\subsection{Memory Management}
\label{sec:memory-management-ch6}

\begin{designnote}
Firmware uses fixed-size buffers and circular/ring buffers for continuous streams. Dynamic memory allocation (\texttt{malloc}/\texttt{new}) is avoided within continuous acquisition loops to reduce heap fragmentation.
\end{designnote}

\subsection{Audio Buffering}
\label{sec:audio-buffering}

\Cref{fig:audio-buffering} shows the audio buffering chain, which helps prevent sample loss during temporary communication delays.

\begin{figure}[H]
\centering
\begin{tikzpicture}[
    font=\small,
    node distance=0mm,
    field/.style={draw, rounded corners=2pt, fill=violet!6, minimum width=4.6cm, minimum height=0.65cm, align=center, line width=0.5pt}
]
\node[field] (mic) {INMP441};
\node[field, below=5mm of mic] (dma) {DMA};
\node[field, below=5mm of dma] (circ) {Circular Buffer};
\node[field, below=5mm of circ] (feat) {Feature Extraction / Packetization};
\node[field, below=5mm of feat] (tx) {Transmission};

\foreach \a/\b in {mic/dma, dma/circ, circ/feat, feat/tx}
  \draw[-{Stealth[length=2mm]}, line width=0.6pt] (\a.south) -- (\b.north);
\end{tikzpicture}
\caption{Audio buffering chain from the INMP441 microphone through DMA and a circular buffer to feature extraction/packetization and transmission, reducing sample loss during transient communication delays.}
\label{fig:audio-buffering}
\end{figure}

\subsection{Firmware Folder Structure}
\label{sec:firmware-folder-structure}

\Cref{lst:firmware-structure} shows the planned firmware repository layout, shared between the Wrist and Audio Module projects.

\begin{figure}[H]
\centering
\begin{lstlisting}[
  basicstyle=\ttfamily\scriptsize,
  frame=single,
  caption={PSYCON firmware repository structure.},
  label={lst:firmware-structure},
  numbers=none
]
firmware/
|-- wrist/
|   |-- sensors/
|   |-- communication/
|   |-- battery/
|   |-- logging/
|   |-- utilities/
|   `-- main.cpp
|-- audio/
|   |-- microphone/
|   |-- communication/
|   |-- battery/
|   |-- logging/
|   `-- main.cpp
`-- shared/
    |-- config.h
    |-- packet.h
    |-- crc.cpp
    `-- timestamp.cpp
\end{lstlisting}
\end{figure}

\subsection{Firmware Validation Checklist}
\label{sec:firmware-validation-checklist}

\begin{requirementbox}
\textbf{Before full integration, the following must be verified:}
\begin{itemize}
    \item Sensor initialization succeeds.
    \item All I\textsuperscript{2}C devices are detected.
    \item No address conflicts.
    \item Stable sampling over extended periods.
    \item No packet corruption.
    \item No memory leaks.
    \item No watchdog resets during normal operation.
    \item Battery warnings trigger correctly.
    \item Graceful shutdown works.
    \item Audio stream remains continuous without excessive buffer overruns.
\end{itemize}
\end{requirementbox}

\subsection{Version Control}
\label{sec:version-control-ch6}

The following are tracked under version control: firmware version; hardware revision; configuration changes; sensor calibration updates; release notes; and test results for each firmware version.

\subsection{Chapter Status}
\label{sec:ch6-status}

\begin{validationbox}
\textbf{Completed for this chapter:}
\begin{itemize}
    \item Firmware architecture
    \item Task scheduling
    \item Sensor acquisition flow
    \item Packet structure
    \item Communication design
    \item Logging strategy
    \item Battery management
    \item Error handling
    \item Memory management
    \item Boot process
    \item Validation checklist
\end{itemize}
\end{validationbox}
